\section{Attenuation with dependent path weights}
\label{sec:path-averages}

The path variance estimate controls sums of edge discrepancies, but
the critical comparison also needs decay of the
products multiplying those sums.  Those products and differences depend
on the same observations.  We therefore average over paths
before using the variance estimate.  Keeping the first half of a path
fixed gives an outward continuation; keeping the second half fixed
requires enumerating the first half backward.  This is why we need
concentration in both directions.

Throughout this section, $b=d-1\ge2$, $\theta=b^{-1/2}$, and
$\omega(t)=\min\{\sqrt t,1\}$.  We first fix $t>0$ and work with the
spins and observations in the planted realization of one marginal.  All fields
below are evaluated at this time, which we usually suppress.  In
particular, $h_{v\to w}$ is the previously chosen measurable function
of the observation history in $T_{v\to w}$ alone.  We use the shifted
estimate from one branch~\eqref{eq:shifted-attenuation}: for an intact
infinite component attached to a spin $\sigma_w$, every real $z$, and
either value $s$ of that spin,
\begin{equation}
 \E\left[a\bigl(z+\psi(h_{v\to w})\bigr)^2
              \mid\sigma_w=s\right]\le1-c_0\omega(t),
 \qquad 0<c_0\le\tfrac12.
 \label{eq:path-shift-input}
\end{equation}
Here $c_0$ depends only on $d$.  The shift may be fixed by conditioning
on other variables as long as the selected component retains its law
conditional on its attaching spin.  This is the only quantitative
information estimate used in the present section.

\subsection{Products in both directions}

For a vertex $z$ and either the empty set or a singleton
$F\subseteq\{u:u\sim z\}$, let $\mathcal P_\ell(z;F)$ consist of the
simple paths $\gamma=(z_0,\ldots,z_\ell)$ that start at $z$ and, when
$\ell\ge1$, have $z_1\notin F$.  Thus
\[
 |\mathcal P_\ell(z;F)|=(d-|F|)b^{\ell-1}\quad(\ell\ge1),
 \qquad \mathcal P_0(z;F)=\{(z)\}.
\]
Define the two squared products and their averages over paths by
\begin{align*}
 R^+_\gamma
  &=\prod_{i=1}^{\ell}a(h_{z_i\to z_{i-1}})^2,
 &G^+_\ell(z;F)
  &=\frac1{|\mathcal P_\ell(z;F)|}
       \sum_{\gamma\in\mathcal P_\ell(z;F)}R^+_\gamma,\\
 R^-_\gamma
  &=\prod_{i=1}^{\ell}a(h_{z_{i-1}\to z_i})^2,
 &G^-_\ell(z;F)
  &=\frac1{|\mathcal P_\ell(z;F)|}
       \sum_{\gamma\in\mathcal P_\ell(z;F)}R^-_\gamma.
\end{align*}
Empty products are one, so $G^+_0=G^-_0=1$.  We call the products
outward and inward according to the coefficient traversal: $R^+$ is
traversed from $z_0$ to $z_\ell$, and $R^-$ from $z_\ell$ to $z_0$.
In each factor the field arrow names the component root and its
excluded neighbor, and points opposite to that coefficient traversal.

We will split the path at an intermediate vertex and average its
continuations inside disjoint descendant components.  Outward
continuation products already use only their descendant observations.
An inward product can also see the observations above the split.  The
next lemma removes the one entering message that causes this dependence.

\begin{lemma}[A continuation below a cut]
\label{lem:cut-continuation}
Let $y$ and $p$ be adjacent vertices, and root the component
$\mathcal T_y=T_{y\to p}$ at $y$.  Write $\operatorname{Ch}(w)$ for the
children of $w$ in this component.  For a descendant path
$y=w_0,w_1,\ldots,w_L$ of length $L\ge1$, set
$H_i=h_{w_i\to w_{i+1}}$ for $0\le i<L$ and define
\begin{align}
 \widetilde H_0
  &=X_y+\sum_{u\in\operatorname{Ch}(y)\setminus\{w_1\}}
                   \psi(h_{u\to y}),\notag\\
 \widetilde H_i
  &=X_{w_i}+\psi(\widetilde H_{i-1})
       +\sum_{u\in\operatorname{Ch}(w_i)\setminus\{w_{i+1}\}}
                   \psi(h_{u\to w_i}),\qquad 1\le i<L.
 \label{eq:inward-cut-recursion}
\end{align}
Each $\widetilde H_i$ is a function only of the observations in
$\mathcal T_y$.  With
\[
 C_{\mathrm{cut}}=\exp\!\left(\frac{4\beta}{1-\theta}\right),
\]
the original and modified inward products obey, almost surely,
\begin{equation}
 \left|\log
   \frac{\prod_{i=0}^{L-1}a(H_i)^2}
        {\prod_{i=0}^{L-1}a(\widetilde H_i)^2}\right|
       \le\frac{4\beta}{1-\theta}.
 \label{eq:path-cut-cost}
\end{equation}
Moreover, conditional on the true spins of this path, both the outward
product and the modified inward product satisfy
\begin{equation}
 \begin{aligned}
 \E\left[\prod_{i=1}^{L}a(h_{w_i\to w_{i-1}})^2
       \,\middle|\,\sigma_{w_0},\ldots,\sigma_{w_L}\right]
       &\le e^{-c_0\omega(t)L},\\
 \E\left[\prod_{i=0}^{L-1}a(\widetilde H_i)^2
       \,\middle|\,\sigma_{w_0},\ldots,\sigma_{w_L}\right]
       &\le e^{-c_0\omega(t)L}.
 \end{aligned}
 \label{eq:cut-continuation-attenuation}
\end{equation}
The constants are independent of $t$, $L$, and the choice of path.
\end{lemma}

\begin{proof}
The recurrence deletes only the contribution from $p$ at $y$ and
propagates that change forward; every field from a side child is unchanged.
All these fields use components contained in $\mathcal T_y$, so the
claimed locality follows by induction.  The modified fields are used
as functions of the descendant observations; no posterior
interpretation for a truncated model is needed.

The cavity identity gives
$H_0-\widetilde H_0=\psi(h_{p\to y})$.  At each subsequent step, only
the predecessor message differs.  Since $|\psi|\le\beta$ and
$|\psi'|\le\theta$, induction yields
\[
 |H_i-\widetilde H_i|\le\beta\theta^i\qquad(0\le i<L).
\]
The bound $|(\log a)'|\le2$ from~\eqref{eq:log-derivative} now gives
\[
 \left|\log
   \frac{\prod_{i=0}^{L-1}a(H_i)^2}
        {\prod_{i=0}^{L-1}a(\widetilde H_i)^2}\right|
 \le4\sum_{i=0}^{L-1}|H_i-\widetilde H_i|
 \le4\beta\sum_{i=0}^{L-1}\theta^i
 \le\frac{4\beta}{1-\theta}.
\]
This proves~\eqref{eq:path-cut-cost}.  In particular, obtaining
component locality costs a multiplicative factor at most
$C_{\mathrm{cut}}$, independent of the length and time.

The outward inequality in~\eqref{eq:cut-continuation-attenuation}
is the conditional path estimate of Lemma~\ref{lem:path-attenuation},
applied to $w_0,\ldots,w_L$.  We prove the modified inward inequality
by integrating intact side components from the far end back toward the cut.

At each target $w_i$, $0\le i<L$, choose one child off the path,
using only the path geometry.  There are $b-1=d-2\ge1$ choices,
including at $w_0$ after deleting its parent and excluding $w_1$.
The selected child components are disjoint and infinite; this also
covers $d=3$ and $L=1$.

Condition on the true path spins.  The selected component observation
histories are independent of each other and of the remaining
observation histories, and their laws depend only on the attaching
path spins.  Indeed, root the broadcast at $w_0$ and group each
selected component's attaching-edge transition, internal transitions,
and Brownian noises.  These groups are mutually independent; once
the path spins are specified, the observations outside the selected
components use none of them.  This factorization passes from finite
cylinder events to the sigma-fields of the full observation histories.
Thus fixing the other observation histories leaves each selected
component with exactly the conditional law used
in~\eqref{eq:path-shift-input}.

Integrate the selected components in the order
$w_{L-1},w_{L-2},\ldots,w_0$.  In~\eqref{eq:inward-cut-recursion},
$\widetilde H_i$ is its selected side message plus a shift fixed by
the other observations.  Every field still present after its factor
is removed has index less than $i$ and does not use that message.
Equation~\eqref{eq:path-shift-input} therefore removes each squared
factor at a cost at most $1-c_0\omega(t)$.  The resulting bound is
$(1-c_0\omega(t))^L\le e^{-c_0\omega(t)L}$, as required.
\end{proof}

\begin{proposition}[Concentration of path averages]
\label{prop:path-average}
Set
\[
 \eta=\min\left\{\frac{c_0}{4},\frac{\log b}{8}\right\},\qquad
 C_0=\max\left\{2,\frac{4\log2}{c_0}\right\},\qquad
 \ell_*(t)=\left\lceil\frac{C_0}{\omega(t)}\right\rceil.
\]
For every $t>0$, every vertex $z$, every allowed first-step restriction
$F$, either orientation $\varepsilon\in\{+,-\}$, and every integer
$\ell\ge\ell_*(t)$,
\begin{equation}
 \P\left(G^\varepsilon_\ell(z;F)
       >C_{\mathrm{cut}}e^{-\eta\omega(t)\ell}\right)
       \le\exp\left(-\frac{b^{\ell/4}}{2b}\right).
 \label{eq:double-exp}
\end{equation}
The constants depend only on $d$ and are uniform in the time, vertex,
first-step restriction, and orientation.
\end{proposition}

\begin{proof}
Root the tree at $z$ and put $k=\lfloor\ell/2\rfloor$ and $L=\ell-k$.
The threshold ensures $k\ge1$.  Let $\mathcal Y_k$ be the vertices at
depth $k$ whose path from $z$ obeys the restriction $F$.  There are
\begin{equation}
 N:=|\mathcal Y_k|=(d-|F|)b^{k-1}\ge b^k
 \label{eq:path-level-size}
\end{equation}
such vertices, and each has exactly $b^L$ descendant continuations
of length $L$.  Since every squared factor lies in $(0,1]$, deleting
the first $k$ factors increases either product.  The average of the
remaining products is consequently an equal-weight average of $N$
continuation averages.

For $y\in\mathcal Y_k$, let $p(y)$ be its parent and
$\mathcal T_y=T_{y\to p(y)}$.  Let $Z_y^+$ be the average of
\[
 \prod_{i=1}^{L}a(h_{w_i\to w_{i-1}})^2
\]
over its $b^L$ continuations $y=w_0,\ldots,w_L$.  Let $Z_y^-$ be
the average of $\prod_{i=0}^{L-1}a(\widetilde H_i)^2$ for the same
continuations, using~\eqref{eq:inward-cut-recursion} separately for
each one.  Both variables lie in $[0,1]$ and are functions only of
the observation histories in $\mathcal T_y$.  Deleting the first
$k$ factors and applying~\eqref{eq:path-cut-cost} gives
\begin{equation}
 G^\varepsilon_\ell(z;F)
       \le\frac{C_{\mathrm{cut}}}{N}
                    \sum_{y\in\mathcal Y_k}Z_y^\varepsilon,
       \qquad \varepsilon\in\{+,-\}.
 \label{eq:path-descendant-average}
\end{equation}
For $+$ the multiplier can be replaced by one.

Averaging~\eqref{eq:cut-continuation-attenuation} first over the
continuation spins conditional on $\sigma_y$, and then over the
continuation choices, yields
\begin{equation}
 \E[Z_y^\varepsilon\mid\sigma_y]
       \le e^{-c_0\omega(t)L},\qquad\varepsilon\in\{+,-\}.
 \label{eq:path-component-mean}
\end{equation}
Now condition on the sigma-field of true cut-level spins
\[
 \mathcal H_k=\sigma(\sigma_y:y\in\mathcal Y_k).
\]
The observations in the disjoint components $\mathcal T_y$ are
independent under this conditioning, with each component law depending
only on its own root spin.  Indeed, after specifying the cut-level
spins, all broadcast transitions and Brownian noises strictly below
the cut, together with the Brownian noise at each cut vertex, can be
sampled independently across these components.  The same statement
holds when $F$ forbids one first neighbor, since that simply omits
one collection of components.  The cut-level spins themselves need
not be independent.  Thus the $Z_y^\varepsilon$ are conditionally
independent variables in $[0,1]$, and their conditional means given
$\mathcal H_k$ satisfy~\eqref{eq:path-component-mean}.

Write $\tau_\ell=e^{-\eta\omega(t)\ell}$.  Since $L\ge\ell/2$,
$\omega(t)\ell\ge C_0$, and $\eta\le c_0/4$, the chosen constants
give
\[
 \E[Z_y^\varepsilon\mid\mathcal H_k]
       \le e^{-c_0\omega(t)\ell/2}\le\tfrac12\tau_\ell.
\]
Apply Hoeffding's bounded-sum inequality
\cite[Theorem~2]{Hoeffding1963} to each conditional product law.
Because the distance from the conditional mean of the average to
$\tau_\ell$ is at least $\tau_\ell/2$, it gives
\begin{align*}
 \P\left(\frac1N\sum_yZ_y^\varepsilon>\tau_\ell
              \,\middle|\,\mathcal H_k\right)
 &\le\exp\left(-2N(\tau_\ell/2)^2\right)\\
 &=\exp\left(-\tfrac12N e^{-2\eta\omega(t)\ell}\right)
 \le\exp\left(-\frac{b^{\ell/4}}{2b}\right).
\end{align*}
For the last inequality, $k\ge\ell/2-1$, $\omega(t)\le1$, and
$2\eta\le(\log b)/4$ imply
\[
 N e^{-2\eta\omega(t)\ell}
       \ge b^{\ell/2-1}e^{-(\log b)\ell/4}
       =b^{\ell/4}/b.
\]
The conditional bound holds almost surely.  Taking
expectations and using~\eqref{eq:path-descendant-average} proves
\eqref{eq:double-exp}.
\end{proof}

Here is the consequence that makes concentration useful for dependent
weights.  Let $\zeta$ be any bounded nonnegative random variable on a
probability space carrying an observation array with the planted
marginal law.  It may also depend on further variables coupled to that
array.  For $\ell\ge\ell_*(t)$, split according to the event in
\eqref{eq:double-exp}; on its complement use the displayed bound for
$G^\varepsilon_\ell$, and on the event use $G^\varepsilon_\ell\le1$.
Then
\begin{equation}
 \E[\zeta G^\varepsilon_\ell(z;F)]
 \le C_{\mathrm{cut}}e^{-\eta\omega(t)\ell}\E\zeta
       +\|\zeta\|_\infty\varepsilon_\ell,
 \qquad
 \varepsilon_\ell:=\exp\left(-\frac{b^{\ell/4}}{2b}\right).
 \label{eq:dependent-average-bound}
\end{equation}
Only the observation marginal enters the event probability.  The
principal term therefore pays the mean of the weight even when the
weight depends on all those observations.

\subsection{Averaging with the discrepancy weights}

Return to the invariant joined law of the two observation copies.
Recall from~\eqref{eq:weighted-path-definitions} that, on a directed
path $v_0,\ldots,v_r$ with $e_i=(v_{i-1},v_i)$,
\[
 P_r=\prod_{i=1}^r a(h_{v_i\to v_{i-1}}),\quad
 Q_r=\E\left[(P_r^2+(P'_r)^2)
                  \left(\sum_{i=1}^rD_{e_i}\right)^2\right],\quad
 L_r=\E[(\Delta h_{v_0})^2P_r^2].
\]
Here $D_{(i,j)}=\log a(h_{j\to i})-\log a(h'_{j\to i})$,
$S=\E D_e^2$, and $V=\E(\Delta h_v)^2$.  Invariance makes $S,V,Q_r,L_r$
independent of the chosen edge, vertex, or directed path of that
length.  The path averages below use that invariance and leave
all observation fields unconditioned.

\begin{proposition}[Weighted path bounds]
\label{prop:weighted-paths}
For every finite $T>0$ there are $c,C>0$, depending only on $d$, and
$C_T<\infty$, depending also on $T$, such that for $0<t\le T$ and
every integer $r\ge2\ell_*(t)$,
\begin{equation}
 \begin{aligned}
 Q_r(t)&\le C e^{-c\omega(t)r}\,rS(t)
       +C_T r^2\exp(-c b^{cr}),\\
 L_r(t)&\le C e^{-c\omega(t)r}\,V(t)
       +C_T\exp(-c b^{cr}).
 \end{aligned}
 \label{eq:weighted-paths}
\end{equation}
At every integer $r\ge1$, without a lower length restriction,
\begin{equation}
 Q_r(t)\le2C_{\mathrm{var}}rS(t),\qquad L_r(t)\le V(t),
 \label{eq:unweighted-bounds}
\end{equation}
where $C_{\mathrm{var}}$ is the constant in
Proposition~\ref{prop:path-variance}.  The hypothesis for
\eqref{eq:weighted-paths} follows in particular from
$r\ge2(C_0+1)/\omega(t)$.
\end{proposition}

\begin{proof}
Fix $T$ and put
\[
 H_T=2T+2d\beta,\qquad
 D_{\max,T}=\frac{2H_T}{1-\theta}.
\]
The site bound~\eqref{eq:bounded-discrepancy}, the component comparison
\eqref{eq:component-difference}, and $|(\log a)'|\le2$ imply
\begin{equation}
 |\Delta h_v|\le H_T,\qquad
 |\Delta h_{j\to i}|
       \le\frac{H_T+\theta H_T}{1-\theta^2}
       =\frac{H_T}{1-\theta},\qquad
 |D_e|\le D_{\max,T}.
 \label{eq:path-discrepancy-bound}
\end{equation}
Dropping $P_r^2,(P'_r)^2\le1$ and applying
Proposition~\ref{prop:path-variance} gives the first bound in
\eqref{eq:unweighted-bounds}; dropping $P_r^2$ gives the second.

For the longer paths, set
\[
 p=\lfloor r/2\rfloor,\qquad q=r-p,\qquad
 U=\sum_{i=1}^pD_{(v_{i-1},v_i)},\qquad
 Z=\sum_{i=p+1}^rD_{(v_{i-1},v_i)}.
\]
Since $(U+Z)^2\le2U^2+2Z^2$,
\begin{equation}
 Q_r\le2\E[(P_r^2+(P'_r)^2)U^2]
       +2\E[(P_r^2+(P'_r)^2)Z^2].
 \label{eq:path-half-split}
\end{equation}
Figure~\ref{fig:half-path-averages} displays the two path averages
that control these terms.

First fix the prefix $(v_0,\ldots,v_p)$ and vary all length-$q$
extensions from $v_p$ that avoid $v_{p-1}$ on their first step.
The pointwise stabilizer of the prefix is transitive on these
extensions: map the successive descendant branches of one extension
to those of the other, and extend independently to the remaining
regular rooted subtrees.  The invariance of the joined law consequently makes
the expectation of the product for each complete path, multiplied
by $U^2$, the same for every extension.  In a fixed joined observation
configuration, $U$ also denotes the same quantity in every term of
this finite average: its indexed directed edges stay fixed.  This
statement places no locality restriction on their fields.  Dropping
the prefix factors from the product gives
\begin{equation}
 \E[P_r^2U^2]
       \le\E[U^2G_q^+(v_p;\{v_{p-1}\})].
 \label{eq:prefix-extension-average}
\end{equation}

For the other term, fix the suffix $(v_p,\ldots,v_r)$ and vary a
reversed prefix $w_0=v_p,w_1,\ldots,w_p$ with $w_1\ne v_{p+1}$.
The resulting complete directed path is
\[
 (w_p,w_{p-1},\ldots,w_1,v_p,v_{p+1},\ldots,v_r).
\]
The pointwise stabilizer of the suffix is transitive on these choices;
it may move the original starting vertex $v_0$, which invariance under
the full automorphism group permits.  The suffix sum $Z$ has the same indexed edges
in every term.  In the direction of the complete path, its prefix
product is exactly
\[
 \prod_{i=1}^p a(h_{w_{i-1}\to w_i})^2,
\]
the inward product along the reversed geometric prefix.  Dropping
the suffix factors therefore yields
\begin{equation}
 \E[P_r^2Z^2]
       \le\E[Z^2G_p^-(v_p;\{v_{p+1}\})].
 \label{eq:suffix-extension-average}
\end{equation}
Each path average is finite.  Thus its use of equivariance
requires only a finite intersection of almost-sure identities for
fixed automorphisms, consistent with the theorem's convention.

\par\addvspace{\intextsep}
\noindent\begin{minipage}{\linewidth}
\makeatletter
\def\@captype{figure}
\makeatother
\centering
\begin{tikzpicture}[
  x=1.55cm,y=1cm,>=Latex,
  every node/.style={font=\small},
  vertex/.style={circle,fill,inner sep=1.5pt},
  fixed/.style={line width=.8pt},
  varied/.style={dashed,line width=.6pt}
]
  \node[vertex,label=left:{$v_0$}] (a) at (0,0) {};
  \node[vertex,label=below:{$v_p$}] (b) at (3,0) {};
  \node[vertex,label=right:{$v_r$}] (c) at (6,0) {};
  \draw[fixed] (a)--(b)
    node[midway,above=3pt] {fixed prefix: $U$};
  \foreach \h in {-.35,0,.35}
    \draw[varied] (b)--(6,\h);
  \draw[->] (3.5,.62)--(5.5,.62);
  \node at (4.5,.98) {average suffix: $G_q^+$};

  \node[vertex,label=left:{$w_p$}] (d) at (0,-2.1) {};
  \node[vertex,label=below:{$v_p$}] (e) at (3,-2.1) {};
  \node[vertex,label=right:{$v_r$}] (f) at (6,-2.1) {};
  \foreach \h in {-.35,0,.35}
    \draw[varied] (e)--(0,-2.1+\h);
  \draw[fixed] (e)--(f)
    node[midway,above=3pt] {fixed suffix: $Z$};
  \draw[->] (2.5,-1.48)--(.5,-1.48);
  \node at (1.5,-1.12) {average reversed prefix: $G_p^-$};
\end{tikzpicture}
\caption{The two averages in the half-path argument.
Solid halves are fixed and dashed halves are varied; each segment
represents a half-path, with one representative pair of endpoints labeled.
Arrows show the direction in which geometric extensions are enumerated.
In the lower row, the prefix is enumerated from $v_p$ toward
$w_p$, while its coefficients are traversed from $w_p$ to $v_p$,
giving the inward average $G_p^-$. The observation fields remain
unconditioned.}
\label{fig:half-path-averages}
\end{minipage}
\par\addvspace{\intextsep}


Suppose $r\ge2\ell_*(t)$.  Then $p,q\ge\ell_*(t)$.  Apply
\eqref{eq:dependent-average-bound} to the right side of
\eqref{eq:prefix-extension-average} with $\zeta=U^2$.  The
unconditional path variance bound and~\eqref{eq:path-discrepancy-bound}
give
$\E U^2\le C_{\mathrm{var}}pS$ and
$\|U^2\|_\infty\le p^2D_{\max,T}^2$.  Hence
\begin{equation}
 \E[P_r^2U^2]
 \le C_{\mathrm{cut}}C_{\mathrm{var}}
          e^{-\eta\omega(t)q}pS+D_{\max,T}^2p^2\varepsilon_q.
 \label{eq:weighted-prefix}
\end{equation}
Applying the same bound to~\eqref{eq:suffix-extension-average}, now
with $\zeta=Z^2$, gives
\begin{equation}
 \E[P_r^2Z^2]
 \le C_{\mathrm{cut}}C_{\mathrm{var}}
          e^{-\eta\omega(t)p}qS+D_{\max,T}^2q^2\varepsilon_p.
 \label{eq:weighted-suffix}
\end{equation}
These applications use the probability, under the joined law, of the
corresponding event determined by one observation marginal.  They do not
require a conditional variance bound or independence of that event
from the discrepancy.  The identical reasoning applies to $(P'_r)^2$
using the second marginal's path averages.

For $r\ge2$, both $p$ and $q$ are at least $r/3$.  Inserting the
four estimates into~\eqref{eq:path-half-split}, and using
$p+q=r$ and $p^2+q^2\le r^2$, proves the explicit inequality
\begin{equation}
 Q_r\le4C_{\mathrm{cut}}C_{\mathrm{var}}
             e^{-\eta\omega(t)r/3}rS
       +4D_{\max,T}^2r^2
             \exp\left(-\frac{b^{r/12}}{2b}\right).
 \label{eq:weighted-path-explicit}
\end{equation}

Finally fix $v_0$ and average the entire directed path over
$\mathcal P_r(v_0;\varnothing)$.  The stabilizer of $v_0$ is transitive on
these paths, and the vertex quantity $\Delta h_{v_0}$ is fixed in
every term.  Invariance gives the exact identity
\[
 L_r=\E[(\Delta h_{v_0})^2G_r^+(v_0;\varnothing)].
\]
For $r\ge\ell_*(t)$, use~\eqref{eq:dependent-average-bound} with
$\zeta=(\Delta h_{v_0})^2$, whose mean is $V$ and whose essential
supremum is at most $H_T^2$.  This yields
\begin{equation}
 L_r\le C_{\mathrm{cut}}e^{-\eta\omega(t)r}V
       +H_T^2\exp\left(-\frac{b^{r/4}}{2b}\right).
 \label{eq:weighted-root-explicit}
\end{equation}
Choosing one sufficiently small positive $c$ in the two explicit
bounds gives~\eqref{eq:weighted-paths}.  Finally,
$\ell_*(t)\le(C_0+1)/\omega(t)$ since $\omega(t)\le1$, which
proves the stated sufficient length condition.
\end{proof}
